\chapter{Background and Related Works}\label{ch:citations}

Android malware detection has attracted significant research attention over the past decade. Researchers have proposed a wide spectrum of techniques—ranging from lightweight static analysis to heavyweight dynamic execution monitoring—each offering different trade-offs between detection accuracy, computational cost, and deployment practicality. This section surveys the key prior works that directly inform VigiDroid's design philosophy, feature engineering choices, model architecture, and on-device deployment strategy.


\section{LinRegDroid: Permission-Based Detection via Multiple Linear Regression}
\label{sec:cite}
Proposed by Narayanan et al.,~\cite{linregdroid} represents one of the early systematic efforts to apply statistical machine learning to Android malware detection using exclusively permission metadata. The central motivation behind this work is that Android permissions requested by an application at install time reveal its underlying behavioral intent. Malicious applications tend to request permissions that are unusual in combination or overly broad relative to their stated utility. LinRegDroid capitalizes on this insight by parsing the AndroidManifest.xml of each APK to extract the requested permissions, encoding them as binary feature flags (1 if requested, 0 if not), and training multiple linear regression classifiers to distinguish between benign and malicious samples. The authors demonstrated that this method is computationally inexpensive, requiring no bytecode extraction or dynamic execution, which makes it ideal for constrained environments.\\

This work directly motivates VigiDroid's exploration of permission-only feature representations as a baseline configuration. The simplicity of permission flags aligns perfectly with VigiDroid core objective: minimizing feature extraction overhead during on-device scanning. VigiDroid acknowledges that while permission signals lack deep expressive power compared to bytecode structural features, they provide a compelling accuracy-to-cost trade-off that makes them viable in a lightweight scanning tier.\\
VigiDroid implements a dedicated LinRegDroid-inspired model configuration within its experimental pipeline. Specifically, the framework includes a multiple linear regression baseline that uses manifest permission flags as the sole feature input. On the temporal holdout of 1,528 APKs from 2022--2023, this configuration reaches 85.5\% accuracy and an F1-score of 0.709, establishing an important lower-bound accuracy reference against which richer feature configurations are compared, while remaining among the fastest detectors at sub-millisecond inference latency.\\

\section{Permission Extraction Framework (MLDP)}
Wang et al.~\cite{mldp} proposed the Malware-relevant Permission Discovery and Pruning (MLDP) framework to address the high dimensionality of Android permission spaces. Rather than feeding hundreds of sparse permission flags directly into a classifier, MLDP ranks permissions by their differential presence rates between malware and benign populations (PRNR), filters candidates by support bounds, and mines association rules over malware-only transactions to retain a compact, discriminative subset---typically fewer than forty permissions. The original paper couples this selection with SVM or tree classifiers trained on the pruned vectors.

VigiDroid adopts the MLDP \emph{feature selection} semantics faithfully but replaces the heavyweight classifiers with a tiny MLP exported to ONNX, consistent with our mobile deployment constraint. The frozen permission set learned on 2020--2021 training data is serialized into the export bundle and replicated byte-for-byte by the Java \texttt{MldpPrunedPermissionExtractor}. This design appears in two cascade tiers: as a standalone Tier~1 model (\texttt{mldp\_pruned\_permission}) and as the stage-one gate inside the MLDP+DEX header cascade (Mode~B).

\section{Detecting Android Malware by Mining Statically Registered Broadcast Receivers}
Mohsen and Ismail presented \cite{broadcast} static analysis approach centered on broadcast receivers as the primary discriminating feature for malware detection. Broadcast receivers are core components that listen for system-wide events (e.g., device boot completion, network status changes, incoming SMS). The authors argue that malicious applications disproportionately register broadcast receivers corresponding to sensitive system events, such as BOOT\_COMPLETED (to achieve persistence after a reboot) or SMS\_RECEIVED (to intercept credentials). Their methodology mines statically declared broadcast receiver actions from the manifest, constructs a presence-absence feature representation, and trains a machine learning classifier. This signal is highly efficient to extract and orthogonal to pure permission-based approaches.\\

This work informs VigiDroid’s inclusion of broadcast receiver actions as an independent feature modality. The insight that receiver patterns carry independent discriminative signals beyond permissions motivates VigiDroid's feature fusion experiments.\\

 VigiDroid's Java pipeline features a dedicated broadcast receiver parser that enumerates all <receiver> elements in the parsed manifest and collects their <intent-filter> action strings. More significantly, this work is synthesized into our best-performing model configuration: Dex+Broadcast Fusion. By combining DEX header metadata with broadcast receiver features, this configuration achieves 96.99\% accuracy and an F1-score of 0.9503 at an exceptionally low CPU time of only 0.9 ms, demonstrating that combining structural bytecode signals with receiver action patterns yields an optimal accuracy-resource frontier.

\section{One-Dimensional Convolutional Neural Networks For Android Malware Detection}
 Hasegawa and Iyatomi \cite{byteCNN} shifted away from structured manifest files to evaluate domain-agnostic, raw byte signals. Their approach trains 1-D Convolutional Neural Networks (CNNs) directly on the raw bytes of an unextracted APK file, testing various sample lengths ($N = 512, 1024, 2048, 4096$) from both the beginning and the end of the file. Crucially, they discovered that analyzing the last $N$ bytes (the tail) of the APK zip archive provides significantly higher classification performance than the first $N$ bytes. Because APK files are fundamentally ZIP archives, the final segments house critical payload structures, directory indices, compiled assets, and occasionally malicious payloads appended to evade typical signature boundaries.\\

 This work serves as the foundational reference for VigiDroids deep learning evaluation tier. It motivates our framework's capacity to bypass structural extraction altogether when a comprehensive file inspection is requested, evaluating whether raw signal processing can out-compete domain-expert feature extraction under rigorous on-device resource constraints.\\ We implemented a self-contained 1-D CNN model that replicates this paradigm. VigiDroid's feature extraction layer includes a tail-byte reader that slices exactly the last 1024 bytes of the raw APK file without extracting its contents. This array is forwarded into a 1-D CNN exported to ONNX. In our multi-tier deployment, this model operates as our Tier 4 heavy-fusion fallback checkpoint, balancing raw sequence-level accuracy against the specialized structural patterns found in earlier tiers.
 \section{SeqMobile: An Efficient Sequence-Based Malware Detection System Using RNN on Mobile Devices}
 Feng et al.\cite{seqmobile} proposed SeqMobile, a pioneering framework for feasible, fully on-device malware detection using Recurrent Neural Networks (RNNs). They extracted both non-sequential features (permissions, intent filters) and sequential features (API call sequences, intent sequences) from the AndroidManifest.xml and the classes.dex binary. To run efficiently on-device, they built specialized feature dictionaries (mapping tokens to unique integers) and implemented native code to extract semantic sequences rapidly. They also optimized the architecture by removing consecutive repetitive calls to reduce sequence sizes by approximately 57\% with minimal accuracy degradation. The trained models were deployed on-device using TensorFlow Lite, comparing dynamic non-quantized and static quantized models.\\
SeqMobile establishes the performance ceiling for on-device sequence processing and directly validates that mobile hardware can handle deep neural architectures if features are properly vectorized. However, SeqMobile's reliance on full API call graph sequencing entails heavy extraction penalties and noticeable runtime delays (2.9s per prediction sequence on a Samsung S10+). This inspired VigiDroid to look for structural shortcuts—such as DEX headers—to capture bytecode characteristics without parsing complete control flow graphs.\\
While VigiDroid replaces heavy sequential RNNs with highly optimized structures like XGBoost and 1-D CNNs to avoid multi-second scanning delays, it adopts SeqMobile's core design choice: performing all feature parsing, vocabulary normalization, and tensor vectorization entirely locally within native-backed components. Furthermore, we utilize their vocabulary dictionary mapping methodology to parse permissions and intent strings into structural indices optimized for ONNX Runtime consumption.

\section{A Lightweight Multi-Source Fast Android Malware Detection Model }
This work \cite{msfdroid} a multi-source fusion methodology utilizing four distinct lightweight base models to perform ensemble learning.
\begin{itemize}
    \item 
Base Model 1 (MLP-H): Extracts and parses structural header bytes from the classes.dex file and passes them to a Multilayer Perceptron.
\item Base Model 2 (MLP-P): Computes the mathematical entropy of the classes.dex file and extracts power spectral density features.
\item Base Model 3 (ASCNN-I): Extracts manifest permissions/intents using a bag-of-words model, shrinking the dimensions via an Adaptive Shrinkage Convolutional Neural Network (ASCNN).
 \item Base Model 4 (ASCNN-C): Concatenates the DEX header features with manifest bag-of-words features through a dual-branch ASCNN architecture.
\end{itemize}
The predictions of these models are fused via an adaptive soft voting ensemble policy and tested live on Android devices.

 This work serves as the architectural bedrock for VigiDroid’s multi-source feature extraction. It highlights that analyzing the structural file metadata of the DEX header provides deep bytecode insights at a fraction of the cost of full disassembling or decompiling. VigiDroid expands on this by introducing temporal evaluation splits and systematically profiling execution resource bounds (battery, memory, CPU time).

 We fully implemented several model variations inspired by this multi-source paradigm. VigiDroid includes a standalone MLP-H configuration that parses the raw classes.dex headers directly from the APK archive. We also implemented the dual-branch feature architecture (Dex+Manifest ASCNN and Dex+Manifest Dual). Most importantly, VigiDroid leverages this work to build a 4-tier cascading scan orchestration system, sequentially deploying these feature extractions in order of computational weight.

\section{ SAMADroid: A Novel 3-Level Hybrid Malware Detection Model}
 Arshad et al. proposed\cite{samadroid} SAMADroid, a hybrid, three-level detection framework designed to combine static and dynamic analysis. The model splits its workload between local and remote environments: lightweight static analysis (permissions, structural manifest components) runs locally on-device to minimize performance overhead, while resource-intensive dynamic analysis (monitoring real-time system calls and network activity) is offloaded to remote host servers.\\

SAMADroid frames the critical trade-off between local processing limits and multi-level enforcement policies. However, its architecture exhibits two clear limitations in modern environments. First, modern Android security models (Android 10 and above) strictly prevent sandboxed applications from launching or tracing other active applications' behaviors, rendering its dynamic capture client obsolete without root privileges. Second, its dependency on active network offloading compromises user privacy and breaks down in offline scenarios. This directly motivated VigiDroid to pursue a fully offline, 100\% on-device cascading pipeline that achieves hybrid-level accuracy without offloading data or violating platform sandbox boundaries.\\

 While VigiDroid completely discards the remote dynamic offloading component due to security and privacy boundaries, it refines SAMADroid's multi-level gating concept. Instead of a local-to-remote split, VigiDroid implements this multi-level concept entirely on-device via a four-tier cascading policy. The initial tiers execute lightweight static checks, and only move to deeper, more complex structural extraction (such as raw byte CNNs) when intermediate classification margins are uncertain.

\section{ DL-Droid: Deep Learning Based Android Malware Detection Using Real Devices}
Alzaylaee et al. \cite{dldroid}created DL-Droid, a framework dedicated to executing deep learning models driven by dynamic features extracted directly from within real physical devices. Their work proved that deep learning classifiers trained on dynamic features (like runtime execution logs and state transitions) can reach exceptionally high performance ceilings and capture zero-day obfuscated malware strains that slip past static signature frameworks.\\

 DL-Droid provides a standard for performance accuracy, but highlights the significant real-world costs of dynamic analysis. Requiring every app to be completely installed and executed for an extended period to catch background hooks is highly impractical for daily on-demand device scanning, consuming immense battery and time. This limitation reinforces VigiDroid’s core hypothesis: lightweight, multi-tiered static structural models can achieve a highly favorable accuracy-resource frontier entirely locally, without requiring app installation or full execution.\\

 VigiDroid uses the performance metrics of DL-Droid as an upper-bound comparison baseline in its thesis evaluation suite. This allows us to explicitly demonstrate how closely VigiDroid's sub-millisecond, on-device cascading static scan modes approach the detection accuracy of comprehensive dynamic frameworks, while completely avoiding execution overhead and device installation requirements.

\section{Drebin: Broad Static Feature Sifting}
Arp et al.~\cite{drebin} introduced a comprehensive static feature set spanning permissions, intents, API calls, and network addresses extracted from decompiled APKs, trained with Support Vector Machines on an OR-encoded vector space. Although Drebin predates modern deep learning, its feature philosophy---cast a wide static net rather than relying on a single manifest signal---informs our XGBoost baseline (\texttt{manifest\_xgb}). VigiDroid's 2,500-dimensional OR vector aggregates manifest and dex-token features in a Drebin-compatible spirit but replaces the SVM with gradient-boosted trees and targets ONNX export for the HEAVY cascade tier. The XGBoost model achieves the highest ROC-AUC (0.991) in our evaluation at the cost of $\sim$100\,ms inference, justifying its placement behind lighter tiers rather than as a first-line detector.


\endinput



